In machine learning, there exists a large number of algorithms, notably neural networks trained with gradient descent, which achieve great performances at prediction tasks on unseen data.
A big question in learning theory is explaining what drives such {\it generalization\/}~\cite{zhang2016understanding,neyshabur2014search}.  
Generalization is the process of applying a model to unseen data, after the algorithm has learned it from some available data.
A number of results have tried to answer this question, coming to the apparent conclusion that the classical theory is not able to exhaustively explain what drives the generalization ability of learning algorithms. In particular, it fails to describe the behavior of the test error of learning algorithms as the complexity of the hypotheses set increases~\cite{belkin2019reconciling,zhang2016understanding,zhang2021understanding, nagarajan2019uniform}.

This thesis seeks to understand whether this is true by testing theoretical results against empirical evidences. We study the setting of a finite hypotheses set of binary decision trees. We look for a bound on the {\it generalization gap\/} as the maximum depth of decision trees increases. The generalization gap is defined as the expected difference between the expected and empirical risk of the selected model, 
\[
\expv_D\left[(R-\hat R_D)(\hat h_D)\right]
\]
where $\expv_D$ denotes the expectation over the disorder of the data. 

We consider a number of candidates coming from existing theoretical results that increasingly account for the structure of the learning problem at hand, in particular the finiteness of decision trees and the noise in the data. Moreover, we test the validity of these results comparing them with the empirical test error curve by carrying out experiments on synthetic data. We are interested in capturing the qualitative behavior of the test error curve. This means that we are not concerned with finding a quantitatively correct bound that lies above the generalization gap, as long as it reproduces its shape. 
By doing so, it is possible to examine how the quantities appearing in the bounds affect their shape, and thereby identify which aspects of the learning problem influence the generalization ability of decision trees. 

The final result of this thesis is a qualitative bound on the generalization gap of decision trees that is empirically validated, providing a theoretical understanding of the behavior of decision trees on unseen data. It also develops framework and toolkit that can be extended to the study of the generalization ability of other learning algorithms, like gradient-based ones.

% in machine learning, comprendere la generalizzazioe degli algoritmi di learning è una domanda centrale. 

% cos'è la generalizzazione, regime overparametrizzato 

% ci sono alcuni lavori che affermano che la teoria statistica classica non sa spiegare bene questa generalizzaizone, soprattutto nel regime overparametrizzato. 

% cosa dicono questi lavori. 

% questo lavoro vuole andare a vedere se questo è vero. 

% per fare ciò, partiamo dal contesto degli alberi decisionali e cerchiamo una funzione maggiorante dell'errore di test degli alberi all'aumentare della profondità massima. ci interessa che questa funzione sia maggiorante in senso asintotico. questo vuol dire che non ci interessa che la curva di tale funzione stia fattualmente sopra alla curva dell'errore di test, bensì che ne riproduca la forma.

% in particolare, andremo a considerare funzioni che considerano sempre di più gli aspetti strutturali del problema di learning in questione, in particolare la finitezza dell'algoritmo e il rumore nei dati. inoltre, testeremo la validità di queste funzioni confrontando la loro forma con quella del'errore di test con gli esperimenti. 

% così facendo, è possibile esaminare come gli oggetti che entrano come input della nostra funzione ne determinano la forma, e dunque quali aspetti del problema di learning influenzano la capacità di generalizzazione degli alberi decisionali. il risultato finale è un bound sul gap (pie l'ha detto bene) degli alberi decisionali che è validato empiricamente e che fornisce una spiegazione teorica della capacità di generalizzazione degli alberi decisionali. 

% il contributo di questo lavoro è 
% qualcosa sugli alberi?
% e fornire un framework che può essere esteso ad altri algoritmi per spiegarne la generalizzazione. 